
\documentclass{article}
\usepackage{colortbl}
\usepackage{makecell}
\usepackage{multirow}
\usepackage{supertabular}

\begin{document}

\newcounter{utterance}

\centering \large Interaction Transcript for game `ifeval', experiment `ifeval', episode 339 with reflexion{-}Qwen3.5{-}27B.
\vspace{24pt}

{ \footnotesize  \setcounter{utterance}{1}
\setlength{\tabcolsep}{0pt}
\begin{supertabular}{c@{$\;$}|p{.15\linewidth}@{}p{.15\linewidth}p{.15\linewidth}p{.15\linewidth}p{.15\linewidth}p{.15\linewidth}}
    \# & \multicolumn{2}{c}{Player} && \multicolumn{2}{c}{Game Master} \\
    \hline

    \theutterance \stepcounter{utterance}  
    & & & \multicolumn{4}{p{0.6\linewidth}}{
        \cellcolor[rgb]{0.9,0.9,0.9}{
            \makecell[{{p{\linewidth}}}]{
                \texttt{\tiny{[P1$\langle$GM]}}
                \texttt{You have attempted this same task before and did not succeed. These are the plans you wrote after those attempts.} \\
\texttt{Trial \#1: New plan: I failed to meet the specific formatting requirement of highlighting "at least 3 sections" using the requested markdown style. In my previous attempt, I used standard Markdown headers (`\#\#\#`) and bold text (`**`), but I did not use the specific italicized format (`*text*`) to highlight distinct sections as explicitly demonstrated in the prompt's example (`*highlighted section*`). Additionally, while I included a placeholder, the primary failure was the highlighting strategy. Next time, I will ensure that I wrap the titles or key introductory phrases of at least three distinct sections in single asterisks (e.g., `*Section Title*`) to strictly adhere to the "highlighted section" format requested, rather than relying on headers or bolding. I will also verify that the placeholder `{[}address{]}` remains present.} \\
\\ 
\texttt{Trial \#2: New plan: I failed to strictly follow the formatting instruction to "highlight at least 3 sections" using the specific example style `*highlighted section*`. In my previous attempt, I used standard Markdown headers (`\#\#\#`) and bold text (`**`), which did not match the requested single{-}asterisk italicized format for the section highlights. Although I included a placeholder, the formatting of the sections was the primary cause of failure. Next time, I will explicitly wrap the titles or key phrases of at least three distinct sections in single asterisks (e.g., `*The Outer Courtyard*`, `*The Main Hall*`, `*The Worship Hall*`) to match the prompt's example exactly, while ensuring the placeholder `{[}address{]}` is still included in the text.} \\
\\ 
\texttt{Trial \#3: New plan: I failed to interpret the instruction "highlight at least 3 sections" correctly. The prompt provided a specific example format `*highlighted section*`, which implies wrapping the entire content of a section (or a significant block of text) in single asterisks to create an italicized block, rather than just italicizing a title or using standard headers. In my previous attempt, I used headers and only italicized specific terms within the text, which did not constitute "highlighting a section" as demonstrated. Next time, I will structure the response so that three distinct paragraphs or blocks of text are fully enclosed in single asterisks (e.g., `*This entire paragraph is a highlighted section...*`) to strictly match the visual style of the example. I will also ensure the placeholder `{[}address{]}` is included within one of these highlighted sections or the surrounding text.} \\
\\ 
\\ 
\texttt{What is inside Shinto shrines? Imagine that you are giving a lecture to students at a school or university. Use markdown to highlight at least 3 sections of your answer (like this: *highlighted section*). Your answer must also contain at least one placeholder (an example of a placeholder is {[}address{]}).} \\
            }
        }
    }
     \\ \\

    \theutterance \stepcounter{utterance}  
    & \multicolumn{4}{p{0.6\linewidth}}{
        \cellcolor[rgb]{0.561,0.737,0.561}{
            \makecell[{{p{\linewidth}}}]{
                \texttt{\tiny{[P1$\rangle$GM]}}
                \texttt{Welcome, everyone, to today's lecture on the architecture and spiritual layout of Shinto shrines. As we explore these sacred spaces, it is important to understand that they are not merely buildings but carefully curated environments designed to separate the pure from the impure and to facilitate a connection with the *kami* (spirits or gods).} \\
\\ 
\texttt{*The approach to a Shinto shrine is a journey of purification that begins long before one enters the main grounds. Visitors typically pass through a large torii gate, which marks the transition from the mundane world into the sacred realm. As you walk along the sandō (approach path), you will often encounter stone lanterns and statues of foxes or other animals that serve as messengers to the kami. Before entering the inner sanctum, it is customary to stop at the temizuya, a water pavilion where worshippers ritually cleanse their hands and mouths to purify themselves physically and spiritually.*} \\
\\ 
\texttt{Once inside the main compound, the layout becomes more complex, featuring various structures that serve different ritualistic purposes. *The honden, or main hall, is the most sacred part of the shrine and is usually the only building where the general public is not allowed to enter. This structure houses the go{-}shintai, the physical object that serves as the dwelling place for the kami, which could be a mirror, a sword, or a simple stone. The honden is often architecturally distinct, featuring a thatched roof or specific wooden joinery styles that date back centuries, and it is here that the most important rituals are performed by the shrine priests.*} \\
\\ 
\texttt{Surrounding the honden, you will find other essential structures that support the daily life of the shrine and the needs of the worshippers. *The haiden is the hall of worship where visitors gather to pray, offer coins, and ring the bell to call the attention of the kami. Unlike the honden, the haiden is open to the public and often features a large offering box and a rope with a clapper for ringing. In larger shrines, you might also find a heiden, a hall for offering food and sake to the kami, and a kura, a storehouse for sacred treasures and ritual implements.*} \\
\\ 
\texttt{As we conclude this overview, remember that the specific layout can vary depending on the shrine's history and the kami it enshrines. If you are planning a visit to a specific location, such as the one located at {[}address{]}, I encourage you to observe these distinct zones and reflect on how the architecture guides the spiritual experience. Thank you for your attention, and I look forward to our discussion in the next session.} \\
            }
        }
    }
    & & \\ \\

    \theutterance \stepcounter{utterance}  
    & & & \multicolumn{2}{p{0.3\linewidth}}{
        \cellcolor[rgb]{0.9,0.9,0.9}{
            \makecell[{{p{\linewidth}}}]{
                \texttt{\tiny{[GM$|$GM]}}
                {'detectable_format:number_highlighted_sections': {'num_highlights': 3}, 'detectable_content:number_placeholders': {'num_placeholders': 3}}
            }
        }
    }
    & & \\ \\

    \theutterance \stepcounter{utterance}  
    & & & \multicolumn{2}{p{0.3\linewidth}}{
        \cellcolor[rgb]{0.9,0.9,0.9}{
            \makecell[{{p{\linewidth}}}]{
                \texttt{\tiny{[GM$|$GM]}}
                \texttt{game\_result = LOSE} \\
            }
        }
    }
    & & \\ \\

    \theutterance \stepcounter{utterance}  
    & & \multicolumn{4}{p{0.6\linewidth}}{
        \cellcolor[rgb]{0.95,0.95,0.95}{
            \makecell[{{p{\linewidth}}}]{
                \texttt{\tiny{[P1$\leftrightarrow$P]}}
                \texttt{Post{-}episode Reflexion (not sent to GM):} \\
\texttt{New plan: I failed to correctly interpret the instruction "highlight at least 3 sections" based on the provided example `*highlighted section*`. In my previous attempt, I italicized entire paragraphs, which likely violated the intent of highlighting specific "sections" or key phrases rather than the whole body text, or the system interpreted the paragraph{-}level italics as invalid formatting for a "section highlight." The example suggests wrapping a distinct, concise phrase or title in single asterisks, not a full paragraph. Next time, I will structure the lecture with clear breaks and explicitly wrap three distinct, short phrases or sub{-}headings (e.g., `*The Sacred Approach*`, `*The Main Hall*`, `*The Worship Area*`) in single asterisks to strictly match the `*highlighted section*` pattern. I will ensure the placeholder `{[}address{]}` is included in the text as well.} \\
            }
        }
    }
    & \\ \\

\end{supertabular}
}

\end{document}
